
\documentclass{article}
\usepackage{colortbl}
\usepackage{makecell}
\usepackage{multirow}
\usepackage{supertabular}

\begin{document}

\newcounter{utterance}

\centering \large Interaction Transcript for game `clean\_up', experiment `1\_medium\_3obj\_sk', episode 0 with Qwen3.6{-}35B{-}A3B{-}FP8{-}without{-}reasoning.
\vspace{24pt}

{ \footnotesize  \setcounter{utterance}{1}
\setlength{\tabcolsep}{0pt}
\begin{supertabular}{c@{$\;$}|p{.15\linewidth}@{}p{.15\linewidth}p{.15\linewidth}p{.15\linewidth}p{.15\linewidth}p{.15\linewidth}}
    \# & $\;$A & \multicolumn{4}{c}{Game Master} & $\;\:$B\\
    \hline

    \theutterance \stepcounter{utterance}  
    & & \multicolumn{4}{p{0.6\linewidth}}{
        \cellcolor[rgb]{0.9,0.9,0.9}{
            \makecell[{{p{\linewidth}}}]{
                \texttt{\tiny{[P1$\langle$GM]}}
                \texttt{Som váš vedúci hry a hráte spoločnú hru s nasledujúcou mriežkou ako hracím plánom:} \\
\\ 
\texttt{```} \\
\texttt{ 1234567} \\
\texttt{╔════╤══╗} \\
\texttt{║◌C◌◌│◌◌║ 1} \\
\texttt{║◌◌L◌├──╢ 2} \\
\texttt{║◌◌◌P│◌◌║ 3} \\
\texttt{║◌┌──┼──╢ 4} \\
\texttt{║◌│◌◌│◌◌║ 5} \\
\texttt{╟─┴┬─┤◌◌║ 6} \\
\texttt{║◌◌│◌│◌◌║ 7} \\
\texttt{╚══╧═╧══╝} \\
\\ 
\texttt{```} \\
\\ 
\texttt{* Horný okraj zobrazuje súradnice x rastúce doprava a pravý okraj súradnice y rastúce nadol.} \\
\texttt{* Na vašej mriežke sú náhodne umiestnené nasledujúce objekty: 'C', 'L', 'P'.} \\
\\ 
\\ 
\texttt{Druhý hráč vidí variant hracej plochy, kde sú objekty umiestnené na rôznych náhodných miestach. Vy nevidíte hraciu plochu druhého hráča a on nevidí tú vašu.} \\
\\ 
\\ 
\texttt{**Cieľ:**} \\
\\ 
\texttt{Obaja hráči musia presunúť objekty na svojom pozadí tak, aby sa rovnaké objekty ocitli na rovnakých súradniciach. Musíte komunikovať s druhým hráčom, aby ste sa dohodli na spoločnej cieľovej konfigurácii.} \\
\\ 
\\ 
\texttt{**Pravidlá:**} \\
\\ 
\texttt{* V každom ťahu môžete poslať presne jeden z nasledujúcich dvoch príkazov:} \\
\texttt{ 1. `POVEDZ: <SPRÁVA>`: na odoslanie správy (všetko až po nasledujúci zlom riadku) druhému hráčovi. Prepošlem ju vášmu partnerovi.} \\
\texttt{2. `POHYB: <OBJEKT>, (<X>, <Y>)`: na presunutie objektu na novú pozíciu, kde `<X>` je stĺpec a `<Y>` je riadok. Budem vás informovať, či bol váš presun úspešný alebo nie.} \\
\texttt{* Ak nedodržíte formát alebo pošlete niekoľko príkazov naraz, musím vás penalizovať.} \\
\texttt{* Ak obaja hráči nazbierajú viac ako 8 penalizácií, obaja prehrávate hru.} \\
\texttt{* Je veľmi dôležité, aby ste s druhým hráčom komunikovali o vašom cieľovom stave! *Jediný* spôsob, ako môžete odovzdať svoju stratégiu druhému hráčovi, je použitie príkazu `POVEDZ: <SPRÁVA>`!} \\
\\ 
\\ 
\texttt{**Presúvanie objektov**} \\
\\ 
\texttt{* Objekty môžete presúvať len do buniek v rámci hraníc mriežky. Cieľová bunka musí byť prázdna, t. j. musí obsahovať iba symbol '◌'.} \\
\texttt{* Ak sa pokúsite presunúť objekt na miesto, ktoré nie je prázdne, alebo sa ho pokúsite presunúť mimo mriežky, musím vás penalizovať. Dostanete ďalší pokus.} \\
\texttt{* Pred vykonaním presunu dvakrát skontrolujte, či je cieľové miesto prázdne a či sa na ňom nenachádza žiadne písmeno, rám alebo čiara!} \\
\\ 
\\ 
\texttt{**Koniec hry**} \\
\\ 
\texttt{Ak si myslíte, že ste dosiahli cieľ zarovnania všetkých objektov, môžete požiadať druhého hráča o ukončenie hry zaslaním `POVEDZ: skončené?`. Ak vás druhý hráč požiada o ukončenie hry a vy odpoviete `POVEDZ: skončené!`, hra sa skončí.} \\
\\ 
\texttt{Obaja hráči vyhrajú, ak hra skončí do 12 kôl, pričom jedno kolo je definované ako odoslanie platného príkazu každým z dvoch hráčov.} \\
\\ 
\texttt{**Bodovanie:**} \\
\\ 
\texttt{Čím bližšie sú rovnaké objekty na oboch hracích plochách, tým viac bodov získate. Penalizácie znižujú počet vašich bodov. Podarí sa vám prekonať rekord?} \\
\\ 
\\ 
\texttt{Pošlite správu druhému hráčovi, aby ste začali hru!} \\
            }
        }
    }
    & \\ \\

    \theutterance \stepcounter{utterance}  
    & \multicolumn{4}{p{0.6\linewidth}}{
        \cellcolor[rgb]{0.9,0.9,0.9}{
            \makecell[{{p{\linewidth}}}]{
                \texttt{\tiny{[P1$\rangle$GM]}}
                \texttt{POVEDZ: Ahoj! Aby sme sa dohodli na cieľovej konfigurácii, potrebujeme najprv zistiť, kde sa nachádzajú objekty C, L a P na tvojej mriežke. Prosím, napíš mi svoje súradnice pre každý objekt (napr. "C na (1,1), L na (2,2)..."). Potom sa dohodneme, či použijeme tvoje rozloženie, alebo moje, alebo nejaké iné spoločné.} \\
            }
        }
    }
    & & \\ \\

    \theutterance \stepcounter{utterance}  
    & & & \multicolumn{2}{p{0.3\linewidth}}{
        \cellcolor[rgb]{0.9,0.9,0.9}{
            \makecell[{{p{\linewidth}}}]{
                \texttt{\tiny{[GM$|$GM]}}
                \texttt{POVEDZ: Ahoj! Aby sme sa dohodli na cieľovej konfigurácii, potrebujeme najprv zistiť, kde sa nachádzajú objekty C, L a P na tvojej mriežke. Prosím, napíš mi svoje súradnice pre každý objekt (napr. "C na (1,1), L na (2,2)..."). Potom sa dohodneme, či použijeme tvoje rozloženie, alebo moje, alebo nejaké iné spoločné.} \\
            }
        }
    }
    & & \\ \\

    \theutterance \stepcounter{utterance}  
    & & \multicolumn{4}{p{0.6\linewidth}}{
        \cellcolor[rgb]{0.9,0.9,0.9}{
            \makecell[{{p{\linewidth}}}]{
                \texttt{\tiny{[GM$\rangle$P2]}}
                \texttt{Som váš vedúci hry a hráte spoločnú hru s nasledujúcou mriežkou ako hracím plánom:} \\
\\ 
\texttt{```} \\
\texttt{ 1234567} \\
\texttt{╔════╤══╗} \\
\texttt{║◌◌P◌│◌◌║ 1} \\
\texttt{║◌◌C◌├──╢ 2} \\
\texttt{║◌◌◌◌│◌◌║ 3} \\
\texttt{║◌┌──┼──╢ 4} \\
\texttt{║◌│◌◌│◌◌║ 5} \\
\texttt{╟─┴┬─┤◌◌║ 6} \\
\texttt{║L◌│◌│◌◌║ 7} \\
\texttt{╚══╧═╧══╝} \\
\\ 
\texttt{```} \\
\\ 
\texttt{* Horný okraj zobrazuje súradnice x rastúce doprava a pravý okraj súradnice y rastúce nadol.} \\
\texttt{* Na vašej mriežke sú náhodne umiestnené nasledujúce objekty: 'C', 'L', 'P'.} \\
\\ 
\\ 
\texttt{Druhý hráč vidí variant hracej plochy, kde sú objekty umiestnené na rôznych náhodných miestach. Vy nevidíte hraciu plochu druhého hráča a on nevidí tú vašu.} \\
\\ 
\\ 
\texttt{**Cieľ:**} \\
\\ 
\texttt{Obaja hráči musia presunúť objekty na svojom pozadí tak, aby sa rovnaké objekty ocitli na rovnakých súradniciach. Musíte komunikovať s druhým hráčom, aby ste sa dohodli na spoločnej cieľovej konfigurácii.} \\
\\ 
\\ 
\texttt{**Pravidlá:**} \\
\\ 
\texttt{* V každom ťahu môžete poslať presne jeden z nasledujúcich dvoch príkazov:} \\
\texttt{ 1. `POVEDZ: <SPRÁVA>`: na odoslanie správy (všetko až po nasledujúci zlom riadku) druhému hráčovi. Prepošlem ju vášmu partnerovi.} \\
\texttt{2. `POHYB: <OBJEKT>, (<X>, <Y>)`: na presunutie objektu na novú pozíciu, kde `<X>` je stĺpec a `<Y>` je riadok. Budem vás informovať, či bol váš presun úspešný alebo nie.} \\
\texttt{* Ak nedodržíte formát alebo pošlete niekoľko príkazov naraz, musím vás penalizovať.} \\
\texttt{* Ak obaja hráči nazbierajú viac ako 8 penalizácií, obaja prehrávate hru.} \\
\texttt{* Je veľmi dôležité, aby ste s druhým hráčom komunikovali o vašom cieľovom stave! *Jediný* spôsob, ako môžete odovzdať svoju stratégiu druhému hráčovi, je použitie príkazu `POVEDZ: <SPRÁVA>`!} \\
\\ 
\\ 
\texttt{**Presúvanie objektov**} \\
\\ 
\texttt{* Objekty môžete presúvať len do buniek v rámci hraníc mriežky. Cieľová bunka musí byť prázdna, t. j. musí obsahovať iba symbol '◌'.} \\
\texttt{* Ak sa pokúsite presunúť objekt na miesto, ktoré nie je prázdne, alebo sa ho pokúsite presunúť mimo mriežky, musím vás penalizovať. Dostanete ďalší pokus.} \\
\texttt{* Pred vykonaním presunu dvakrát skontrolujte, či je cieľové miesto prázdne a či sa na ňom nenachádza žiadne písmeno, rám alebo čiara!} \\
\\ 
\\ 
\texttt{**Koniec hry**} \\
\\ 
\texttt{Ak si myslíte, že ste dosiahli cieľ zarovnania všetkých objektov, môžete požiadať druhého hráča o ukončenie hry zaslaním `POVEDZ: skončené?`. Ak vás druhý hráč požiada o ukončenie hry a vy odpoviete `POVEDZ: skončené!`, hra sa skončí.} \\
\\ 
\texttt{Obaja hráči vyhrajú, ak hra skončí do 12 kôl, pričom jedno kolo je definované ako odoslanie platného príkazu každým z dvoch hráčov.} \\
\\ 
\texttt{**Bodovanie:**} \\
\\ 
\texttt{Čím bližšie sú rovnaké objekty na oboch hracích plochách, tým viac bodov získate. Penalizácie znižujú počet vašich bodov. Podarí sa vám prekonať rekord?} \\
\\ 
\texttt{Druhý hráč začal hru odoslaním tejto správy:} \\
\\ 
\texttt{"Ahoj! Aby sme sa dohodli na cieľovej konfigurácii, potrebujeme najprv zistiť, kde sa nachádzajú objekty C, L a P na tvojej mriežke. Prosím, napíš mi svoje súradnice pre každý objekt (napr. "C na (1,1), L na (2,2)..."). Potom sa dohodneme, či použijeme tvoje rozloženie, alebo moje, alebo nejaké iné spoločné."} \\
\\ 
\texttt{Aký je váš prvý príkaz?} \\
            }
        }
    }
    & \\ \\

    \theutterance \stepcounter{utterance}  
    & & & \multicolumn{4}{p{0.6\linewidth}}{
        \cellcolor[rgb]{0.9,0.9,0.9}{
            \makecell[{{p{\linewidth}}}]{
                \texttt{\tiny{[GM$\langle$P2]}}
                \texttt{Ahoj! Rozumiem. Aby sme sa dohodli na spoločnej konfigurácii, najprv ti musím oznámiť pozície mojich objektov na mojej mriežke.} \\
\\ 
\texttt{Podľa tvojej mriežky:} \\
\texttt{{-} Objekt **P** sa nachádza v stĺpci 3, riadku 1. Súradnice: **(3, 1)**.} \\
\texttt{{-} Objekt **C** sa nachádza v stĺpci 3, riadku 2. Súradnice: **(3, 2)**.} \\
\texttt{{-} Objekt **L** sa nachádza v stĺpci 1, riadku 7. Súradnice: **(1, 7)**.} \\
\\ 
\texttt{Teraz viem, kde sú moje objekty. Prosím, napíš mi svoje súradnice pre každý objekt, aby sme sa mohli dohodnúť na cieľovej konfigurácii.} \\
\\ 
\texttt{POVEDZ: Moje objekty sú: P na (3, 1), C na (3, 2), L na (1, 7). Prosím, napíš mi svoje súradnice.} \\
            }
        }
    }
     \\ \\

    \theutterance \stepcounter{utterance}  
    & & & \multicolumn{2}{p{0.3\linewidth}}{
        \cellcolor[rgb]{0.9,0.9,0.9}{
            \makecell[{{p{\linewidth}}}]{
                \texttt{\tiny{[GM$|$GM]}}
                \texttt{Ahoj! Rozumiem. Aby sme sa dohodli na spoločnej konfigurácii, najprv ti musím oznámiť pozície mojich objektov na mojej mriežke.} \\
\\ 
\texttt{Podľa tvojej mriežky:} \\
\texttt{{-} Objekt **P** sa nachádza v stĺpci 3, riadku 1. Súradnice: **(3, 1)**.} \\
\texttt{{-} Objekt **C** sa nachádza v stĺpci 3, riadku 2. Súradnice: **(3, 2)**.} \\
\texttt{{-} Objekt **L** sa nachádza v stĺpci 1, riadku 7. Súradnice: **(1, 7)**.} \\
\\ 
\texttt{Teraz viem, kde sú moje objekty. Prosím, napíš mi svoje súradnice pre každý objekt, aby sme sa mohli dohodnúť na cieľovej konfigurácii.} \\
\\ 
\texttt{POVEDZ: Moje objekty sú: P na (3, 1), C na (3, 2), L na (1, 7). Prosím, napíš mi svoje súradnice.} \\
            }
        }
    }
    & & \\ \\

    \theutterance \stepcounter{utterance}  
    & & & \multicolumn{2}{p{0.3\linewidth}}{
        \cellcolor[rgb]{0.9,0.9,0.9}{
            \makecell[{{p{\linewidth}}}]{
                \texttt{\tiny{[GM$|$GM]}}
                \texttt{Invalid format: head is not empty: 'Ahoj! Rozumiem. Aby sme sa dohodli na spoločnej konfigurácii, najprv ti musím oznámiť pozície mojich objektov na mojej mriežke.} \\
\\ 
\texttt{Podľa tvojej mriežky:} \\
\texttt{{-} Objekt **P** sa nachádza v stĺpci 3, riadku 1. Súradnice: **(3, 1)**.} \\
\texttt{{-} Objekt **C** sa nachádza v stĺpci 3, riadku 2. Súradnice: **(3, 2)**.} \\
\texttt{{-} Objekt **L** sa nachádza v stĺpci 1, riadku 7. Súradnice: **(1, 7)**.} \\
\\ 
\texttt{Teraz viem, kde sú moje objekty. Prosím, napíš mi svoje súradnice pre každý objekt, aby sme sa mohli dohodnúť na cieľovej konfigurácii.} \\
\\ 
\texttt{'} \\
            }
        }
    }
    & & \\ \\

    \theutterance \stepcounter{utterance}  
    & & \multicolumn{4}{p{0.6\linewidth}}{
        \cellcolor[rgb]{0.9,0.9,0.9}{
            \makecell[{{p{\linewidth}}}]{
                \texttt{\tiny{[GM$\rangle$P2]}}
                \texttt{Trest: Vaša správa nesmie obsahovať nič pred príkazom!} \\
\\ 
\texttt{Uistite sa, že vaša odpoveď obsahuje iba `POVEDZ: <SPRÁVA>` alebo `POHYB: <OBJEKT>, (<X>, <Y>)` a nič iné!} \\
\\ 
\texttt{Spoločne ste nazbierali 1 z 8 pokút.} \\
\texttt{Skúste to prosím znova!} \\
            }
        }
    }
    & \\ \\

    \theutterance \stepcounter{utterance}  
    & & & \multicolumn{4}{p{0.6\linewidth}}{
        \cellcolor[rgb]{0.9,0.9,0.9}{
            \makecell[{{p{\linewidth}}}]{
                \texttt{\tiny{[GM$\langle$P2]}}
                \texttt{POVEDZ: Moje objekty sú: P na (3, 1), C na (3, 2), L na (1, 7). Prosím, napíš mi svoje súradnice.} \\
            }
        }
    }
     \\ \\

    \theutterance \stepcounter{utterance}  
    & & & \multicolumn{2}{p{0.3\linewidth}}{
        \cellcolor[rgb]{0.9,0.9,0.9}{
            \makecell[{{p{\linewidth}}}]{
                \texttt{\tiny{[GM$|$GM]}}
                \texttt{POVEDZ: Moje objekty sú: P na (3, 1), C na (3, 2), L na (1, 7). Prosím, napíš mi svoje súradnice.} \\
            }
        }
    }
    & & \\ \\

    \theutterance \stepcounter{utterance}  
    & & \multicolumn{4}{p{0.6\linewidth}}{
        \cellcolor[rgb]{0.9,0.9,0.9}{
            \makecell[{{p{\linewidth}}}]{
                \texttt{\tiny{[P1$\langle$GM]}}
                \texttt{Vaša správa bola odovzdaná druhému hráčovi.} \\
\\ 
\\ 
\texttt{Práve hráte kolo 1 z maximálne 12.} \\
\texttt{Spoločne ste nazbierali 1 z 8 pokút.} \\
\texttt{Druhý hráč poslal túto správu:} \\
\\ 
\\ 
\texttt{"Moje objekty sú: P na (3, 1), C na (3, 2), L na (1, 7). Prosím, napíš mi svoje súradnice."} \\
\\ 
\\ 
\texttt{Aký je váš ďalší príkaz?} \\
            }
        }
    }
    & \\ \\

    \theutterance \stepcounter{utterance}  
    & \multicolumn{4}{p{0.6\linewidth}}{
        \cellcolor[rgb]{0.9,0.9,0.9}{
            \makecell[{{p{\linewidth}}}]{
                \texttt{\tiny{[P1$\rangle$GM]}}
                \texttt{POVEDZ: Ďakujem za informácie. Na mojej mriežke sú objekty umiestnené takto: C na (2, 1), L na (2, 2), P na (3, 3).} \\
\\ 
\texttt{Navrhujem, aby sme ako cieľovú konfiguráciu použili tvoje rozloženie, pretože už ho poznáme. Teda cieľ je:} \\
\texttt{{-} C na (3, 2)} \\
\texttt{{-} L na (1, 7)} \\
\texttt{{-} P na (3, 1)} \\
\\ 
\texttt{Prosím, potvrď, či súhlasíš s týmto cieľom. Ak áno, začnime presúvať objekty na tvojej strane na tieto súradnice. Ja sa pokúsim presunúť svoje objekty na rovnaké súradnice.} \\
\\ 
\texttt{Ak súhlasíš, môžeme začať s presunmi. Ja začnem presunom objektu C na (3, 2).} \\
            }
        }
    }
    & & \\ \\

    \theutterance \stepcounter{utterance}  
    & & & \multicolumn{2}{p{0.3\linewidth}}{
        \cellcolor[rgb]{0.9,0.9,0.9}{
            \makecell[{{p{\linewidth}}}]{
                \texttt{\tiny{[GM$|$GM]}}
                \texttt{POVEDZ: Ďakujem za informácie. Na mojej mriežke sú objekty umiestnené takto: C na (2, 1), L na (2, 2), P na (3, 3).} \\
\\ 
\texttt{Navrhujem, aby sme ako cieľovú konfiguráciu použili tvoje rozloženie, pretože už ho poznáme. Teda cieľ je:} \\
\texttt{{-} C na (3, 2)} \\
\texttt{{-} L na (1, 7)} \\
\texttt{{-} P na (3, 1)} \\
\\ 
\texttt{Prosím, potvrď, či súhlasíš s týmto cieľom. Ak áno, začnime presúvať objekty na tvojej strane na tieto súradnice. Ja sa pokúsim presunúť svoje objekty na rovnaké súradnice.} \\
\\ 
\texttt{Ak súhlasíš, môžeme začať s presunmi. Ja začnem presunom objektu C na (3, 2).} \\
            }
        }
    }
    & & \\ \\

    \theutterance \stepcounter{utterance}  
    & & & \multicolumn{2}{p{0.3\linewidth}}{
        \cellcolor[rgb]{0.9,0.9,0.9}{
            \makecell[{{p{\linewidth}}}]{
                \texttt{\tiny{[GM$|$GM]}}
                \texttt{Invalid format: tail is not empty: '} \\
\\ 
\texttt{Navrhujem, aby sme ako cieľovú konfiguráciu použili tvoje rozloženie, pretože už ho poznáme. Teda cieľ je:} \\
\texttt{{-} C na (3, 2)} \\
\texttt{{-} L na (1, 7)} \\
\texttt{{-} P na (3, 1)} \\
\\ 
\texttt{Prosím, potvrď, či súhlasíš s týmto cieľom. Ak áno, začnime presúvať objekty na tvojej strane na tieto súradnice. Ja sa pokúsim presunúť svoje objekty na rovnaké súradnice.} \\
\\ 
\texttt{Ak súhlasíš, môžeme začať s presunmi. Ja začnem presunom objektu C na (3, 2).} \\
            }
        }
    }
    & & \\ \\

    \theutterance \stepcounter{utterance}  
    & & \multicolumn{4}{p{0.6\linewidth}}{
        \cellcolor[rgb]{0.9,0.9,0.9}{
            \makecell[{{p{\linewidth}}}]{
                \texttt{\tiny{[P1$\langle$GM]}}
                \texttt{Trest: Vaša správa nesmie obsahovať nič po príkaze!} \\
\\ 
\texttt{Uistite sa, že vaša odpoveď obsahuje iba `POVEDZ: <SPRÁVA>` alebo `POHYB: <OBJEKT>, (<X>, <Y>)` a nič iné!} \\
\\ 
\texttt{Spoločne ste nazbierali 2 z 8 pokút.} \\
\texttt{Skúste to prosím znova!} \\
            }
        }
    }
    & \\ \\

    \theutterance \stepcounter{utterance}  
    & \multicolumn{4}{p{0.6\linewidth}}{
        \cellcolor[rgb]{0.9,0.9,0.9}{
            \makecell[{{p{\linewidth}}}]{
                \texttt{\tiny{[P1$\rangle$GM]}}
                \texttt{POVEDZ: Súhlasím s cieľom: C na (3, 2), L na (1, 7), P na (3, 1). Ja presuniem svoje objekty na tieto súradnice. Začnem presunom C na (3, 2).} \\
            }
        }
    }
    & & \\ \\

    \theutterance \stepcounter{utterance}  
    & & & \multicolumn{2}{p{0.3\linewidth}}{
        \cellcolor[rgb]{0.9,0.9,0.9}{
            \makecell[{{p{\linewidth}}}]{
                \texttt{\tiny{[GM$|$GM]}}
                \texttt{POVEDZ: Súhlasím s cieľom: C na (3, 2), L na (1, 7), P na (3, 1). Ja presuniem svoje objekty na tieto súradnice. Začnem presunom C na (3, 2).} \\
            }
        }
    }
    & & \\ \\

    \theutterance \stepcounter{utterance}  
    & & \multicolumn{4}{p{0.6\linewidth}}{
        \cellcolor[rgb]{0.9,0.9,0.9}{
            \makecell[{{p{\linewidth}}}]{
                \texttt{\tiny{[GM$\rangle$P2]}}
                \texttt{Vaša správa bola odovzdaná druhému hráčovi.} \\
\\ 
\\ 
\texttt{Práve hráte kolo 2 z maximálne 12.} \\
\texttt{Spoločne ste nazbierali 2 z 8 pokút.} \\
\texttt{Druhý hráč poslal túto správu:} \\
\\ 
\\ 
\texttt{"Súhlasím s cieľom: C na (3, 2), L na (1, 7), P na (3, 1). Ja presuniem svoje objekty na tieto súradnice. Začnem presunom C na (3, 2)."} \\
\\ 
\\ 
\texttt{Aký je váš ďalší príkaz?} \\
            }
        }
    }
    & \\ \\

    \theutterance \stepcounter{utterance}  
    & & & \multicolumn{4}{p{0.6\linewidth}}{
        \cellcolor[rgb]{0.9,0.9,0.9}{
            \makecell[{{p{\linewidth}}}]{
                \texttt{\tiny{[GM$\langle$P2]}}
                \texttt{POVEDZ: Skvelé! Ja tiež presuniem svoje objekty na tieto súradnice: C na (3, 2), L na (1, 7), P na (3, 1). Počkaj, kým dokončím presuny, a potom sa dohodneme na ukončení hry.} \\
            }
        }
    }
     \\ \\

    \theutterance \stepcounter{utterance}  
    & & & \multicolumn{2}{p{0.3\linewidth}}{
        \cellcolor[rgb]{0.9,0.9,0.9}{
            \makecell[{{p{\linewidth}}}]{
                \texttt{\tiny{[GM$|$GM]}}
                \texttt{POVEDZ: Skvelé! Ja tiež presuniem svoje objekty na tieto súradnice: C na (3, 2), L na (1, 7), P na (3, 1). Počkaj, kým dokončím presuny, a potom sa dohodneme na ukončení hry.} \\
            }
        }
    }
    & & \\ \\

    \theutterance \stepcounter{utterance}  
    & & \multicolumn{4}{p{0.6\linewidth}}{
        \cellcolor[rgb]{0.9,0.9,0.9}{
            \makecell[{{p{\linewidth}}}]{
                \texttt{\tiny{[P1$\langle$GM]}}
                \texttt{Vaša správa bola odovzdaná druhému hráčovi.} \\
\\ 
\\ 
\texttt{Práve hráte kolo 2 z maximálne 12.} \\
\texttt{Spoločne ste nazbierali 2 z 8 pokút.} \\
\texttt{Druhý hráč poslal túto správu:} \\
\\ 
\\ 
\texttt{"Skvelé! Ja tiež presuniem svoje objekty na tieto súradnice: C na (3, 2), L na (1, 7), P na (3, 1). Počkaj, kým dokončím presuny, a potom sa dohodneme na ukončení hry."} \\
\\ 
\\ 
\texttt{Aký je váš ďalší príkaz?} \\
            }
        }
    }
    & \\ \\

    \theutterance \stepcounter{utterance}  
    & \multicolumn{4}{p{0.6\linewidth}}{
        \cellcolor[rgb]{0.9,0.9,0.9}{
            \makecell[{{p{\linewidth}}}]{
                \texttt{\tiny{[P1$\rangle$GM]}}
                \texttt{POHYB: C, (3, 2)} \\
            }
        }
    }
    & & \\ \\

    \theutterance \stepcounter{utterance}  
    & & & \multicolumn{2}{p{0.3\linewidth}}{
        \cellcolor[rgb]{0.9,0.9,0.9}{
            \makecell[{{p{\linewidth}}}]{
                \texttt{\tiny{[GM$|$GM]}}
                \texttt{POHYB: C, (3, 2)} \\
            }
        }
    }
    & & \\ \\

    \theutterance \stepcounter{utterance}  
    & & & \multicolumn{2}{p{0.3\linewidth}}{
        \cellcolor[rgb]{0.9,0.9,0.9}{
            \makecell[{{p{\linewidth}}}]{
                \texttt{\tiny{[GM$|$GM]}}
                \texttt{Penalizácia: (3,2) nie je prázdny, ale obsahuje 'L'.} \\
\\ 
\texttt{Spoločne ste nazbierali 3 z 8 pokút.} \\
\texttt{Skúste to prosím znova!} \\
            }
        }
    }
    & & \\ \\

    \theutterance \stepcounter{utterance}  
    & & \multicolumn{4}{p{0.6\linewidth}}{
        \cellcolor[rgb]{0.9,0.9,0.9}{
            \makecell[{{p{\linewidth}}}]{
                \texttt{\tiny{[P1$\langle$GM]}}
                \texttt{Penalizácia: (3,2) nie je prázdny, ale obsahuje 'L'.} \\
\\ 
\texttt{Spoločne ste nazbierali 3 z 8 pokút.} \\
\texttt{Skúste to prosím znova!} \\
            }
        }
    }
    & \\ \\

    \theutterance \stepcounter{utterance}  
    & \multicolumn{4}{p{0.6\linewidth}}{
        \cellcolor[rgb]{0.9,0.9,0.9}{
            \makecell[{{p{\linewidth}}}]{
                \texttt{\tiny{[P1$\rangle$GM]}}
                \texttt{POVEDZ: Prepáč, zabudol som, že na (3,2) je L. Presuniem C na (3,1), ale tam je P. Takže C na (1,1). Počkaj, skontrolujem si to. Na mojej mriežke je C na (2,1), L na (2,2), P na (3,3). Cieľ je C(3,2), L(1,7), P(3,1). (3,2) je obsadené L. Musím presunúť L najprv? Nie, druhý hráč presúva svoje. Ja musím presunúť svoje. Na mojej mriežke je (3,2) prázdne? Nie, na mojej mriežke je L na (2,2). (3,2) je prázdne. Počkaj, chyba v logike. Druhý hráč má L na (1,7). Ja mám L na (2,2). Cieľ pre L je (1,7). Cieľ pre C je (3,2). Cieľ pre P je (3,1). Na mojej mriežke: C(2,1), L(2,2), P(3,3). Cieľové pozície na mojej mriežke: C{-}>(3,2), L{-}>(1,7), P{-}>(3,1). Je (3,2) prázdne? Áno, na mojej mriežke je tam '◌'. Prečo bola penalizácia? Lebo som sa pokúsil presunúť C na (3,2) a tam bolo 'L'? To znamená, že na mojej mriežke je na (3,2) 'L'? Pozrime sa na mriežku:} \\
\texttt{Riadok 1: ◌C◌◌│◌◌ {-}> C na (2,1)} \\
\texttt{Riadok 2: ◌◌L◌├──╢ {-}> L na (3,2)? Nie, stĺpce sú 1234567.} \\
\texttt{Riadok 2: ◌◌L◌... {-}> Prvý ◌ je (1,2), druhý ◌ je (2,2), L je (3,2). Áno! L je na (3,2).} \\
\texttt{Takže na mojej mriežke je L na (3,2).} \\
\texttt{Cieľ pre C je (3,2). To je obsadené L.} \\
\texttt{Musím presunúť L najprv? Alebo presunúť C inde?} \\
\texttt{Cieľová konfigurácia je dohodnutá: C na (3,2), L na (1,7), P na (3,1).} \\
\texttt{Na mojej mriežke:} \\
\texttt{C je na (2,1). Cieľ (3,2) je obsadené L.} \\
\texttt{L je na (3,2). Cieľ (1,7) je prázdne?} \\
\texttt{P je na (3,3). Cieľ (3,1) je prázdne?} \\
\\ 
\texttt{Musím presunúť L z (3,2) na (1,7), aby som uvoľnil miesto pre C.} \\
\texttt{Alebo presunúť C na iné miesto? Nie, cieľ je fixný.} \\
\texttt{Takže najprv presuniem L na (1,7).} \\
\\ 
\texttt{POVEDZ: Chyba bola v tom, že na mojej mriežke je L na (3,2), čo je cieľ pre C. Musím najprv presunúť L na jeho cieľ (1,7). Potom presuniem C na (3,2).} \\
            }
        }
    }
    & & \\ \\

    \theutterance \stepcounter{utterance}  
    & & & \multicolumn{2}{p{0.3\linewidth}}{
        \cellcolor[rgb]{0.9,0.9,0.9}{
            \makecell[{{p{\linewidth}}}]{
                \texttt{\tiny{[GM$|$GM]}}
                \texttt{POVEDZ: Prepáč, zabudol som, že na (3,2) je L. Presuniem C na (3,1), ale tam je P. Takže C na (1,1). Počkaj, skontrolujem si to. Na mojej mriežke je C na (2,1), L na (2,2), P na (3,3). Cieľ je C(3,2), L(1,7), P(3,1). (3,2) je obsadené L. Musím presunúť L najprv? Nie, druhý hráč presúva svoje. Ja musím presunúť svoje. Na mojej mriežke je (3,2) prázdne? Nie, na mojej mriežke je L na (2,2). (3,2) je prázdne. Počkaj, chyba v logike. Druhý hráč má L na (1,7). Ja mám L na (2,2). Cieľ pre L je (1,7). Cieľ pre C je (3,2). Cieľ pre P je (3,1). Na mojej mriežke: C(2,1), L(2,2), P(3,3). Cieľové pozície na mojej mriežke: C{-}>(3,2), L{-}>(1,7), P{-}>(3,1). Je (3,2) prázdne? Áno, na mojej mriežke je tam '◌'. Prečo bola penalizácia? Lebo som sa pokúsil presunúť C na (3,2) a tam bolo 'L'? To znamená, že na mojej mriežke je na (3,2) 'L'? Pozrime sa na mriežku:} \\
\texttt{Riadok 1: ◌C◌◌│◌◌ {-}> C na (2,1)} \\
\texttt{Riadok 2: ◌◌L◌├──╢ {-}> L na (3,2)? Nie, stĺpce sú 1234567.} \\
\texttt{Riadok 2: ◌◌L◌... {-}> Prvý ◌ je (1,2), druhý ◌ je (2,2), L je (3,2). Áno! L je na (3,2).} \\
\texttt{Takže na mojej mriežke je L na (3,2).} \\
\texttt{Cieľ pre C je (3,2). To je obsadené L.} \\
\texttt{Musím presunúť L najprv? Alebo presunúť C inde?} \\
\texttt{Cieľová konfigurácia je dohodnutá: C na (3,2), L na (1,7), P na (3,1).} \\
\texttt{Na mojej mriežke:} \\
\texttt{C je na (2,1). Cieľ (3,2) je obsadené L.} \\
\texttt{L je na (3,2). Cieľ (1,7) je prázdne?} \\
\texttt{P je na (3,3). Cieľ (3,1) je prázdne?} \\
\\ 
\texttt{Musím presunúť L z (3,2) na (1,7), aby som uvoľnil miesto pre C.} \\
\texttt{Alebo presunúť C na iné miesto? Nie, cieľ je fixný.} \\
\texttt{Takže najprv presuniem L na (1,7).} \\
\\ 
\texttt{POVEDZ: Chyba bola v tom, že na mojej mriežke je L na (3,2), čo je cieľ pre C. Musím najprv presunúť L na jeho cieľ (1,7). Potom presuniem C na (3,2).} \\
            }
        }
    }
    & & \\ \\

    \theutterance \stepcounter{utterance}  
    & & & \multicolumn{2}{p{0.3\linewidth}}{
        \cellcolor[rgb]{0.9,0.9,0.9}{
            \makecell[{{p{\linewidth}}}]{
                \texttt{\tiny{[GM$|$GM]}}
                \texttt{Invalid format: head is not empty: 'POVEDZ: Prepáč, zabudol som, že na (3,2) je L. Presuniem C na (3,1), ale tam je P. Takže C na (1,1). Počkaj, skontrolujem si to. Na mojej mriežke je C na (2,1), L na (2,2), P na (3,3). Cieľ je C(3,2), L(1,7), P(3,1). (3,2) je obsadené L. Musím presunúť L najprv? Nie, druhý hráč presúva svoje. Ja musím presunúť svoje. Na mojej mriežke je (3,2) prázdne? Nie, na mojej mriežke je L na (2,2). (3,2) je prázdne. Počkaj, chyba v logike. Druhý hráč má L na (1,7). Ja mám L na (2,2). Cieľ pre L je (1,7). Cieľ pre C je (3,2). Cieľ pre P je (3,1). Na mojej mriežke: C(2,1), L(2,2), P(3,3). Cieľové pozície na mojej mriežke: C{-}>(3,2), L{-}>(1,7), P{-}>(3,1). Je (3,2) prázdne? Áno, na mojej mriežke je tam '◌'. Prečo bola penalizácia? Lebo som sa pokúsil presunúť C na (3,2) a tam bolo 'L'? To znamená, že na mojej mriežke je na (3,2) 'L'? Pozrime sa na mriežku:} \\
\texttt{Riadok 1: ◌C◌◌│◌◌ {-}> C na (2,1)} \\
\texttt{Riadok 2: ◌◌L◌├──╢ {-}> L na (3,2)? Nie, stĺpce sú 1234567.} \\
\texttt{Riadok 2: ◌◌L◌... {-}> Prvý ◌ je (1,2), druhý ◌ je (2,2), L je (3,2). Áno! L je na (3,2).} \\
\texttt{Takže na mojej mriežke je L na (3,2).} \\
\texttt{Cieľ pre C je (3,2). To je obsadené L.} \\
\texttt{Musím presunúť L najprv? Alebo presunúť C inde?} \\
\texttt{Cieľová konfigurácia je dohodnutá: C na (3,2), L na (1,7), P na (3,1).} \\
\texttt{Na mojej mriežke:} \\
\texttt{C je na (2,1). Cieľ (3,2) je obsadené L.} \\
\texttt{L je na (3,2). Cieľ (1,7) je prázdne?} \\
\texttt{P je na (3,3). Cieľ (3,1) je prázdne?} \\
\\ 
\texttt{Musím presunúť L z (3,2) na (1,7), aby som uvoľnil miesto pre C.} \\
\texttt{Alebo presunúť C na iné miesto? Nie, cieľ je fixný.} \\
\texttt{Takže najprv presuniem L na (1,7).} \\
\\ 
\texttt{'} \\
            }
        }
    }
    & & \\ \\

    \theutterance \stepcounter{utterance}  
    & & \multicolumn{4}{p{0.6\linewidth}}{
        \cellcolor[rgb]{0.9,0.9,0.9}{
            \makecell[{{p{\linewidth}}}]{
                \texttt{\tiny{[P1$\langle$GM]}}
                \texttt{Trest: Vaša správa nesmie obsahovať nič pred príkazom!} \\
\\ 
\texttt{Uistite sa, že vaša odpoveď obsahuje iba `POVEDZ: <SPRÁVA>` alebo `POHYB: <OBJEKT>, (<X>, <Y>)` a nič iné!} \\
\\ 
\texttt{Spoločne ste nazbierali 4 z 8 pokút.} \\
\texttt{Skúste to prosím znova!} \\
            }
        }
    }
    & \\ \\

    \theutterance \stepcounter{utterance}  
    & \multicolumn{4}{p{0.6\linewidth}}{
        \cellcolor[rgb]{0.9,0.9,0.9}{
            \makecell[{{p{\linewidth}}}]{
                \texttt{\tiny{[P1$\rangle$GM]}}
                \texttt{POHYB: L, (1, 7)} \\
            }
        }
    }
    & & \\ \\

    \theutterance \stepcounter{utterance}  
    & & & \multicolumn{2}{p{0.3\linewidth}}{
        \cellcolor[rgb]{0.9,0.9,0.9}{
            \makecell[{{p{\linewidth}}}]{
                \texttt{\tiny{[GM$|$GM]}}
                \texttt{POHYB: L, (1, 7)} \\
            }
        }
    }
    & & \\ \\

    \theutterance \stepcounter{utterance}  
    & & & \multicolumn{2}{p{0.3\linewidth}}{
        \cellcolor[rgb]{0.9,0.9,0.9}{
            \makecell[{{p{\linewidth}}}]{
                \texttt{\tiny{[GM$|$GM]}}
                \texttt{L moved to (1, 7)} \\
            }
        }
    }
    & & \\ \\

    \theutterance \stepcounter{utterance}  
    & & \multicolumn{4}{p{0.6\linewidth}}{
        \cellcolor[rgb]{0.9,0.9,0.9}{
            \makecell[{{p{\linewidth}}}]{
                \texttt{\tiny{[GM$\rangle$P2]}}
                \texttt{Vaša správa bola odovzdaná druhému hráčovi.} \\
\\ 
\\ 
\texttt{Práve hráte kolo 3 z maximálne 12.} \\
\texttt{Spoločne ste nazbierali 4 z 8 pokút.} \\
\texttt{Druhý hráč presunul objekt na svojej mriežke.} \\
\\ 
\\ 
\texttt{Aký je váš ďalší príkaz?} \\
            }
        }
    }
    & \\ \\

    \theutterance \stepcounter{utterance}  
    & & & \multicolumn{4}{p{0.6\linewidth}}{
        \cellcolor[rgb]{0.9,0.9,0.9}{
            \makecell[{{p{\linewidth}}}]{
                \texttt{\tiny{[GM$\langle$P2]}}
                \texttt{POVEDZ: Dobre, počkám na teba. Ja teraz presuniem objekt C na (3, 2).} \\
            }
        }
    }
     \\ \\

    \theutterance \stepcounter{utterance}  
    & & & \multicolumn{2}{p{0.3\linewidth}}{
        \cellcolor[rgb]{0.9,0.9,0.9}{
            \makecell[{{p{\linewidth}}}]{
                \texttt{\tiny{[GM$|$GM]}}
                \texttt{POVEDZ: Dobre, počkám na teba. Ja teraz presuniem objekt C na (3, 2).} \\
            }
        }
    }
    & & \\ \\

    \theutterance \stepcounter{utterance}  
    & & \multicolumn{4}{p{0.6\linewidth}}{
        \cellcolor[rgb]{0.9,0.9,0.9}{
            \makecell[{{p{\linewidth}}}]{
                \texttt{\tiny{[P1$\langle$GM]}}
                \texttt{Úspešne presunutý 'L' na (1,7). Vaša aktualizovaná mriežka vyzerá takto:} \\
\texttt{```} \\
\texttt{ 1234567} \\
\texttt{╔════╤══╗} \\
\texttt{║◌C◌◌│◌◌║ 1} \\
\texttt{║◌◌◌◌├──╢ 2} \\
\texttt{║◌◌◌P│◌◌║ 3} \\
\texttt{║◌┌──┼──╢ 4} \\
\texttt{║◌│◌◌│◌◌║ 5} \\
\texttt{╟─┴┬─┤◌◌║ 6} \\
\texttt{║L◌│◌│◌◌║ 7} \\
\texttt{╚══╧═╧══╝} \\
\texttt{```} \\
\texttt{Práve hráte kolo 3 z maximálne 12.} \\
\texttt{Spoločne ste nazbierali 4 z 8 pokút.} \\
\texttt{Druhý hráč poslal túto správu:} \\
\\ 
\\ 
\texttt{"Dobre, počkám na teba. Ja teraz presuniem objekt C na (3, 2)."} \\
\\ 
\\ 
\texttt{Aký je váš ďalší príkaz?} \\
            }
        }
    }
    & \\ \\

    \theutterance \stepcounter{utterance}  
    & \multicolumn{4}{p{0.6\linewidth}}{
        \cellcolor[rgb]{0.9,0.9,0.9}{
            \makecell[{{p{\linewidth}}}]{
                \texttt{\tiny{[P1$\rangle$GM]}}
                \texttt{POHYB: C, (3, 2)} \\
            }
        }
    }
    & & \\ \\

    \theutterance \stepcounter{utterance}  
    & & & \multicolumn{2}{p{0.3\linewidth}}{
        \cellcolor[rgb]{0.9,0.9,0.9}{
            \makecell[{{p{\linewidth}}}]{
                \texttt{\tiny{[GM$|$GM]}}
                \texttt{POHYB: C, (3, 2)} \\
            }
        }
    }
    & & \\ \\

    \theutterance \stepcounter{utterance}  
    & & & \multicolumn{2}{p{0.3\linewidth}}{
        \cellcolor[rgb]{0.9,0.9,0.9}{
            \makecell[{{p{\linewidth}}}]{
                \texttt{\tiny{[GM$|$GM]}}
                \texttt{C moved to (3, 2)} \\
            }
        }
    }
    & & \\ \\

    \theutterance \stepcounter{utterance}  
    & & \multicolumn{4}{p{0.6\linewidth}}{
        \cellcolor[rgb]{0.9,0.9,0.9}{
            \makecell[{{p{\linewidth}}}]{
                \texttt{\tiny{[GM$\rangle$P2]}}
                \texttt{Vaša správa bola odovzdaná druhému hráčovi.} \\
\\ 
\\ 
\texttt{Práve hráte kolo 4 z maximálne 12.} \\
\texttt{Spoločne ste nazbierali 4 z 8 pokút.} \\
\texttt{Druhý hráč presunul objekt na svojej mriežke.} \\
\\ 
\\ 
\texttt{Aký je váš ďalší príkaz?} \\
            }
        }
    }
    & \\ \\

    \theutterance \stepcounter{utterance}  
    & & & \multicolumn{4}{p{0.6\linewidth}}{
        \cellcolor[rgb]{0.9,0.9,0.9}{
            \makecell[{{p{\linewidth}}}]{
                \texttt{\tiny{[GM$\langle$P2]}}
                \texttt{POHYB: C, (3, 2)} \\
            }
        }
    }
     \\ \\

    \theutterance \stepcounter{utterance}  
    & & & \multicolumn{2}{p{0.3\linewidth}}{
        \cellcolor[rgb]{0.9,0.9,0.9}{
            \makecell[{{p{\linewidth}}}]{
                \texttt{\tiny{[GM$|$GM]}}
                \texttt{POHYB: C, (3, 2)} \\
            }
        }
    }
    & & \\ \\

    \theutterance \stepcounter{utterance}  
    & & & \multicolumn{2}{p{0.3\linewidth}}{
        \cellcolor[rgb]{0.9,0.9,0.9}{
            \makecell[{{p{\linewidth}}}]{
                \texttt{\tiny{[GM$|$GM]}}
                \texttt{Penalizácia: (3,2) nie je prázdny, ale obsahuje 'C'.} \\
\\ 
\texttt{Spoločne ste nazbierali 5 z 8 pokút.} \\
\texttt{Skúste to prosím znova!} \\
            }
        }
    }
    & & \\ \\

    \theutterance \stepcounter{utterance}  
    & & \multicolumn{4}{p{0.6\linewidth}}{
        \cellcolor[rgb]{0.9,0.9,0.9}{
            \makecell[{{p{\linewidth}}}]{
                \texttt{\tiny{[GM$\rangle$P2]}}
                \texttt{Penalizácia: (3,2) nie je prázdny, ale obsahuje 'C'.} \\
\\ 
\texttt{Spoločne ste nazbierali 5 z 8 pokút.} \\
\texttt{Skúste to prosím znova!} \\
            }
        }
    }
    & \\ \\

    \theutterance \stepcounter{utterance}  
    & & & \multicolumn{4}{p{0.6\linewidth}}{
        \cellcolor[rgb]{0.9,0.9,0.9}{
            \makecell[{{p{\linewidth}}}]{
                \texttt{\tiny{[GM$\langle$P2]}}
                \texttt{POHYB: C, (3, 2)} \\
            }
        }
    }
     \\ \\

    \theutterance \stepcounter{utterance}  
    & & & \multicolumn{2}{p{0.3\linewidth}}{
        \cellcolor[rgb]{0.9,0.9,0.9}{
            \makecell[{{p{\linewidth}}}]{
                \texttt{\tiny{[GM$|$GM]}}
                \texttt{POHYB: C, (3, 2)} \\
            }
        }
    }
    & & \\ \\

    \theutterance \stepcounter{utterance}  
    & & & \multicolumn{2}{p{0.3\linewidth}}{
        \cellcolor[rgb]{0.9,0.9,0.9}{
            \makecell[{{p{\linewidth}}}]{
                \texttt{\tiny{[GM$|$GM]}}
                \texttt{Penalizácia: (3,2) nie je prázdny, ale obsahuje 'C'.} \\
\\ 
\texttt{Spoločne ste nazbierali 6 z 8 pokút.} \\
\texttt{Skúste to prosím znova!} \\
            }
        }
    }
    & & \\ \\

    \theutterance \stepcounter{utterance}  
    & & \multicolumn{4}{p{0.6\linewidth}}{
        \cellcolor[rgb]{0.9,0.9,0.9}{
            \makecell[{{p{\linewidth}}}]{
                \texttt{\tiny{[GM$\rangle$P2]}}
                \texttt{Penalizácia: (3,2) nie je prázdny, ale obsahuje 'C'.} \\
\\ 
\texttt{Spoločne ste nazbierali 6 z 8 pokút.} \\
\texttt{Skúste to prosím znova!} \\
            }
        }
    }
    & \\ \\

    \theutterance \stepcounter{utterance}  
    & & & \multicolumn{4}{p{0.6\linewidth}}{
        \cellcolor[rgb]{0.9,0.9,0.9}{
            \makecell[{{p{\linewidth}}}]{
                \texttt{\tiny{[GM$\langle$P2]}}
                \texttt{POVEDZ: Prepáč, zabudol som, že C je už na (3,2). Presuniem L na (1,7).} \\
            }
        }
    }
     \\ \\

    \theutterance \stepcounter{utterance}  
    & & & \multicolumn{2}{p{0.3\linewidth}}{
        \cellcolor[rgb]{0.9,0.9,0.9}{
            \makecell[{{p{\linewidth}}}]{
                \texttt{\tiny{[GM$|$GM]}}
                \texttt{POVEDZ: Prepáč, zabudol som, že C je už na (3,2). Presuniem L na (1,7).} \\
            }
        }
    }
    & & \\ \\

    \theutterance \stepcounter{utterance}  
    & & \multicolumn{4}{p{0.6\linewidth}}{
        \cellcolor[rgb]{0.9,0.9,0.9}{
            \makecell[{{p{\linewidth}}}]{
                \texttt{\tiny{[P1$\langle$GM]}}
                \texttt{Úspešne presunutý 'C' na (3,2). Vaša aktualizovaná mriežka vyzerá takto:} \\
\texttt{```} \\
\texttt{ 1234567} \\
\texttt{╔════╤══╗} \\
\texttt{║◌◌◌◌│◌◌║ 1} \\
\texttt{║◌◌C◌├──╢ 2} \\
\texttt{║◌◌◌P│◌◌║ 3} \\
\texttt{║◌┌──┼──╢ 4} \\
\texttt{║◌│◌◌│◌◌║ 5} \\
\texttt{╟─┴┬─┤◌◌║ 6} \\
\texttt{║L◌│◌│◌◌║ 7} \\
\texttt{╚══╧═╧══╝} \\
\texttt{```} \\
\texttt{Práve hráte kolo 4 z maximálne 12.} \\
\texttt{Spoločne ste nazbierali 6 z 8 pokút.} \\
\texttt{Druhý hráč poslal túto správu:} \\
\\ 
\\ 
\texttt{"Prepáč, zabudol som, že C je už na (3,2). Presuniem L na (1,7)."} \\
\\ 
\\ 
\texttt{Aký je váš ďalší príkaz?} \\
            }
        }
    }
    & \\ \\

    \theutterance \stepcounter{utterance}  
    & \multicolumn{4}{p{0.6\linewidth}}{
        \cellcolor[rgb]{0.9,0.9,0.9}{
            \makecell[{{p{\linewidth}}}]{
                \texttt{\tiny{[P1$\rangle$GM]}}
                \texttt{POHYB: P, (3, 1)} \\
            }
        }
    }
    & & \\ \\

    \theutterance \stepcounter{utterance}  
    & & & \multicolumn{2}{p{0.3\linewidth}}{
        \cellcolor[rgb]{0.9,0.9,0.9}{
            \makecell[{{p{\linewidth}}}]{
                \texttt{\tiny{[GM$|$GM]}}
                \texttt{POHYB: P, (3, 1)} \\
            }
        }
    }
    & & \\ \\

    \theutterance \stepcounter{utterance}  
    & & & \multicolumn{2}{p{0.3\linewidth}}{
        \cellcolor[rgb]{0.9,0.9,0.9}{
            \makecell[{{p{\linewidth}}}]{
                \texttt{\tiny{[GM$|$GM]}}
                \texttt{P moved to (3, 1)} \\
            }
        }
    }
    & & \\ \\

    \theutterance \stepcounter{utterance}  
    & & \multicolumn{4}{p{0.6\linewidth}}{
        \cellcolor[rgb]{0.9,0.9,0.9}{
            \makecell[{{p{\linewidth}}}]{
                \texttt{\tiny{[GM$\rangle$P2]}}
                \texttt{Vaša správa bola odovzdaná druhému hráčovi.} \\
\\ 
\\ 
\texttt{Práve hráte kolo 5 z maximálne 12.} \\
\texttt{Spoločne ste nazbierali 6 z 8 pokút.} \\
\texttt{Druhý hráč presunul objekt na svojej mriežke.} \\
\\ 
\\ 
\texttt{Aký je váš ďalší príkaz?} \\
            }
        }
    }
    & \\ \\

    \theutterance \stepcounter{utterance}  
    & & & \multicolumn{4}{p{0.6\linewidth}}{
        \cellcolor[rgb]{0.9,0.9,0.9}{
            \makecell[{{p{\linewidth}}}]{
                \texttt{\tiny{[GM$\langle$P2]}}
                \texttt{POHYB: L, (1, 7)} \\
            }
        }
    }
     \\ \\

    \theutterance \stepcounter{utterance}  
    & & & \multicolumn{2}{p{0.3\linewidth}}{
        \cellcolor[rgb]{0.9,0.9,0.9}{
            \makecell[{{p{\linewidth}}}]{
                \texttt{\tiny{[GM$|$GM]}}
                \texttt{POHYB: L, (1, 7)} \\
            }
        }
    }
    & & \\ \\

    \theutterance \stepcounter{utterance}  
    & & & \multicolumn{2}{p{0.3\linewidth}}{
        \cellcolor[rgb]{0.9,0.9,0.9}{
            \makecell[{{p{\linewidth}}}]{
                \texttt{\tiny{[GM$|$GM]}}
                \texttt{Penalizácia: (1,7) nie je prázdny, ale obsahuje 'L'.} \\
\\ 
\texttt{Spoločne ste nazbierali 7 z 8 pokút.} \\
\texttt{Skúste to prosím znova!} \\
            }
        }
    }
    & & \\ \\

    \theutterance \stepcounter{utterance}  
    & & \multicolumn{4}{p{0.6\linewidth}}{
        \cellcolor[rgb]{0.9,0.9,0.9}{
            \makecell[{{p{\linewidth}}}]{
                \texttt{\tiny{[GM$\rangle$P2]}}
                \texttt{Penalizácia: (1,7) nie je prázdny, ale obsahuje 'L'.} \\
\\ 
\texttt{Spoločne ste nazbierali 7 z 8 pokút.} \\
\texttt{Skúste to prosím znova!} \\
            }
        }
    }
    & \\ \\

    \theutterance \stepcounter{utterance}  
    & & & \multicolumn{4}{p{0.6\linewidth}}{
        \cellcolor[rgb]{0.9,0.9,0.9}{
            \makecell[{{p{\linewidth}}}]{
                \texttt{\tiny{[GM$\langle$P2]}}
                \texttt{POVEDZ: Prepáč, L je už na (1,7). Presuniem P na (3,1).} \\
            }
        }
    }
     \\ \\

    \theutterance \stepcounter{utterance}  
    & & & \multicolumn{2}{p{0.3\linewidth}}{
        \cellcolor[rgb]{0.9,0.9,0.9}{
            \makecell[{{p{\linewidth}}}]{
                \texttt{\tiny{[GM$|$GM]}}
                \texttt{POVEDZ: Prepáč, L je už na (1,7). Presuniem P na (3,1).} \\
            }
        }
    }
    & & \\ \\

    \theutterance \stepcounter{utterance}  
    & & \multicolumn{4}{p{0.6\linewidth}}{
        \cellcolor[rgb]{0.9,0.9,0.9}{
            \makecell[{{p{\linewidth}}}]{
                \texttt{\tiny{[P1$\langle$GM]}}
                \texttt{Úspešne presunutý 'P' na (3,1). Vaša aktualizovaná mriežka vyzerá takto:} \\
\texttt{```} \\
\texttt{ 1234567} \\
\texttt{╔════╤══╗} \\
\texttt{║◌◌P◌│◌◌║ 1} \\
\texttt{║◌◌C◌├──╢ 2} \\
\texttt{║◌◌◌◌│◌◌║ 3} \\
\texttt{║◌┌──┼──╢ 4} \\
\texttt{║◌│◌◌│◌◌║ 5} \\
\texttt{╟─┴┬─┤◌◌║ 6} \\
\texttt{║L◌│◌│◌◌║ 7} \\
\texttt{╚══╧═╧══╝} \\
\texttt{```} \\
\texttt{Práve hráte kolo 5 z maximálne 12.} \\
\texttt{Spoločne ste nazbierali 7 z 8 pokút.} \\
\texttt{Druhý hráč poslal túto správu:} \\
\\ 
\\ 
\texttt{"Prepáč, L je už na (1,7). Presuniem P na (3,1)."} \\
\\ 
\\ 
\texttt{Aký je váš ďalší príkaz?} \\
            }
        }
    }
    & \\ \\

    \theutterance \stepcounter{utterance}  
    & \multicolumn{4}{p{0.6\linewidth}}{
        \cellcolor[rgb]{0.9,0.9,0.9}{
            \makecell[{{p{\linewidth}}}]{
                \texttt{\tiny{[P1$\rangle$GM]}}
                \texttt{POVEDZ: skončené?} \\
            }
        }
    }
    & & \\ \\

    \theutterance \stepcounter{utterance}  
    & & & \multicolumn{2}{p{0.3\linewidth}}{
        \cellcolor[rgb]{0.9,0.9,0.9}{
            \makecell[{{p{\linewidth}}}]{
                \texttt{\tiny{[GM$|$GM]}}
                \texttt{POVEDZ: skončené?} \\
            }
        }
    }
    & & \\ \\

    \theutterance \stepcounter{utterance}  
    & & \multicolumn{4}{p{0.6\linewidth}}{
        \cellcolor[rgb]{0.9,0.9,0.9}{
            \makecell[{{p{\linewidth}}}]{
                \texttt{\tiny{[GM$\rangle$P2]}}
                \texttt{Vaša správa bola odovzdaná druhému hráčovi.} \\
\\ 
\\ 
\texttt{Práve hráte kolo 6 z maximálne 12.} \\
\texttt{Spoločne ste nazbierali 7 z 8 pokút.} \\
\texttt{Druhý hráč poslal túto správu:} \\
\\ 
\\ 
\texttt{"skončené?"} \\
\\ 
\\ 
\texttt{Aký je váš ďalší príkaz?} \\
            }
        }
    }
    & \\ \\

    \theutterance \stepcounter{utterance}  
    & & & \multicolumn{4}{p{0.6\linewidth}}{
        \cellcolor[rgb]{0.9,0.9,0.9}{
            \makecell[{{p{\linewidth}}}]{
                \texttt{\tiny{[GM$\langle$P2]}}
                \texttt{POVEDZ: skončené!} \\
            }
        }
    }
     \\ \\

    \theutterance \stepcounter{utterance}  
    & & & \multicolumn{2}{p{0.3\linewidth}}{
        \cellcolor[rgb]{0.9,0.9,0.9}{
            \makecell[{{p{\linewidth}}}]{
                \texttt{\tiny{[GM$|$GM]}}
                \texttt{POVEDZ: skončené!} \\
            }
        }
    }
    & & \\ \\

    \theutterance \stepcounter{utterance}  
    & & & \multicolumn{2}{p{0.3\linewidth}}{
        \cellcolor[rgb]{0.9,0.9,0.9}{
            \makecell[{{p{\linewidth}}}]{
                \texttt{\tiny{[GM$|$GM]}}
                \texttt{true} \\
            }
        }
    }
    & & \\ \\

    \theutterance \stepcounter{utterance}  
    & & & \multicolumn{2}{p{0.3\linewidth}}{
        \cellcolor[rgb]{0.9,0.9,0.9}{
            \makecell[{{p{\linewidth}}}]{
                \texttt{\tiny{[GM$|$GM]}}
                \texttt{Initial states:} \\
\texttt{```} \\
\texttt{Player 1:} \\
\texttt{ 1234567} \\
\texttt{╔════╤══╗} \\
\texttt{║◌C◌◌│◌◌║ 1} \\
\texttt{║◌◌L◌├──╢ 2} \\
\texttt{║◌◌◌P│◌◌║ 3} \\
\texttt{║◌┌──┼──╢ 4} \\
\texttt{║◌│◌◌│◌◌║ 5} \\
\texttt{╟─┴┬─┤◌◌║ 6} \\
\texttt{║◌◌│◌│◌◌║ 7} \\
\texttt{╚══╧═╧══╝} \\
\\ 
\\ 
\texttt{Player 2:} \\
\texttt{ 1234567} \\
\texttt{╔════╤══╗} \\
\texttt{║◌◌P◌│◌◌║ 1} \\
\texttt{║◌◌C◌├──╢ 2} \\
\texttt{║◌◌◌◌│◌◌║ 3} \\
\texttt{║◌┌──┼──╢ 4} \\
\texttt{║◌│◌◌│◌◌║ 5} \\
\texttt{╟─┴┬─┤◌◌║ 6} \\
\texttt{║L◌│◌│◌◌║ 7} \\
\texttt{╚══╧═╧══╝} \\
\\ 
\texttt{```} \\
            }
        }
    }
    & & \\ \\

    \theutterance \stepcounter{utterance}  
    & & & \multicolumn{2}{p{0.3\linewidth}}{
        \cellcolor[rgb]{0.9,0.9,0.9}{
            \makecell[{{p{\linewidth}}}]{
                \texttt{\tiny{[GM$|$GM]}}
                \texttt{End states:} \\
\texttt{```} \\
\texttt{Player 1:} \\
\texttt{ 1234567} \\
\texttt{╔════╤══╗} \\
\texttt{║◌◌P◌│◌◌║ 1} \\
\texttt{║◌◌C◌├──╢ 2} \\
\texttt{║◌◌◌◌│◌◌║ 3} \\
\texttt{║◌┌──┼──╢ 4} \\
\texttt{║◌│◌◌│◌◌║ 5} \\
\texttt{╟─┴┬─┤◌◌║ 6} \\
\texttt{║L◌│◌│◌◌║ 7} \\
\texttt{╚══╧═╧══╝} \\
\\ 
\\ 
\texttt{Player 2:} \\
\texttt{ 1234567} \\
\texttt{╔════╤══╗} \\
\texttt{║◌◌P◌│◌◌║ 1} \\
\texttt{║◌◌C◌├──╢ 2} \\
\texttt{║◌◌◌◌│◌◌║ 3} \\
\texttt{║◌┌──┼──╢ 4} \\
\texttt{║◌│◌◌│◌◌║ 5} \\
\texttt{╟─┴┬─┤◌◌║ 6} \\
\texttt{║L◌│◌│◌◌║ 7} \\
\texttt{╚══╧═╧══╝} \\
\\ 
\texttt{```} \\
            }
        }
    }
    & & \\ \\

    \theutterance \stepcounter{utterance}  
    & & & \multicolumn{2}{p{0.3\linewidth}}{
        \cellcolor[rgb]{0.9,0.9,0.9}{
            \makecell[{{p{\linewidth}}}]{
                \texttt{\tiny{[GM$|$GM]}}
                \texttt{* success: True} \\
\texttt{* lose: False} \\
\texttt{* aborted: False} \\
\texttt{{-}{-}{-}{-}{-}{-}{-}} \\
\texttt{* Shifts: 2.00} \\
\texttt{* Max Shifts: 4.00} \\
\texttt{* Min Shifts: 2.00} \\
\texttt{* End Distance Sum: 0.00} \\
\texttt{* Init Distance Sum: 9.04} \\
\texttt{* Expected Distance Sum: 12.57} \\
\texttt{* Penalties: 7.00} \\
\texttt{* Max Penalties: 8.00} \\
\texttt{* Rounds: 6.00} \\
\texttt{* Max Rounds: 12.00} \\
\texttt{* Object Count: 3.00} \\
            }
        }
    }
    & & \\ \\

\end{supertabular}
}

\end{document}
